\chapter{Syllables}

All syllables must have exactly one vowel.
Neither syllabic consonants nor vowel hiatuses are allowed.
Subsequent vowels must separated by at least one intervening consonant.

In native items, \phonem{d, d͡ʒ, g, ɣ, l, ŋ, ɾ, w, z} do not appear word\-/%
initially, and \phonem{b, d͡ʒ, d, g, ɣ} do not appear in codas.
Additionally, complex onsets do not exist, and complex codas are subject to
restrictions.
Only \phonem{jt, nt͡ʃ, nt, ɾk, ɾq, ɾs, ɾt} are allowed as complex codas.
Thus, the maximal syllable structure in native items is (C)V(C)(C), with
examples of (C)V syllables in~(\ref{ex:(c)v-syllables}), (C)VC in~(\ref{%
ex:(c)vc-syllables}), and (C)VCC in~(\ref{ex:(c)vcc-syllables}).

\pex
  \a\label{ex:(c)v-syllables}
    \vtop{\halign{
      \phonem{#}\hfil&\:\transl{#}\hfil\cr
      ɛ.ki&two\cr
      bu&this\cr
    }}
  \a\label{ex:(c)vc-syllables}
    \vtop{\halign{
      \phonem{#}\hfil&\:\transl{#}\hfil\cr
      ɑl.tɯ&six\cr
      jɯɾ&song\cr
    }}
  \a\label{ex:(c)vcc-syllables}
    \vtop{\halign{
      \phonem{#}\hfil&\:\transl{#}\hfil\cr
      ɑnt&oath\cr
      qɯɾq&forty\cr
    }}
\xe

Loans may admit syllables of different types from those describe above, such as
those with complex onsets, other types of complex codas, or less restricted
distributions of consonants.

\ex
  \vtop{\halign{
    \phonem{#}\hfil&\:\transl{#}\hfil\cr
    hɑd͡ʒ&pilgrimage\cr
    dɔst&friend\cr
    pɾɔ.gɾɑm.mɑ&computer program\cr
  }}
\xe
